\documentclass[11pt,a4paper]{article}
\usepackage[T1]{fontenc}
\usepackage{lmodern,amsmath,amssymb,amsthm,mathtools,booktabs,microtype}
\usepackage[margin=27mm]{geometry}
\usepackage[colorlinks=true,linkcolor=blue,citecolor=blue,urlcolor=blue]{hyperref}
\newtheorem{theorem}{Theorem}[section]
\newtheorem{lemma}[theorem]{Lemma}
\newtheorem{proposition}[theorem]{Proposition}
\newcommand{\Q}{\mathbb Q}\newcommand{\Z}{\mathbb Z}
\setlength{\emergencystretch}{2em}
\title{A quartic in 257 variables with an effective parameter}
\author{}\date{16 November 2026}
\begin{document}\maketitle

\begin{abstract}
An affine centering and a Laurent elimination give a completely
specified quartic family in 257 rational variables, including the
input. The construction has 229 quadratic equations and 3679 nonzero
quartic monomials. It retains the effective sufficient parameter
$N=9682$ and the arithmetic premises of the explicit source, without
an additional $abc$ assumption. At this parameter its integer
coefficient height is exactly $149^{19366}$. The transformation is
an input-preserving bijection on the full rational zero sets.
\end{abstract}

\section{Statement and arithmetic scope}

Write $C=149^N$. The source is the explicit 258-coordinate system
of~\cite{Laurent,Explicit}, in its balanced integer quadratic table
\begin{center}\small
\path{computation/quartic/laurent_recovery_v38/family258/balanced-quadratic.json}.
\end{center}
Its SHA-256 digest is
\begin{center}\small\ttfamily
ad8204d383210613590b3daa8031dc43113d50b0d26e84d067971cac5f37fe4c.
\end{center}
We retain exactly its elliptic-group and chart assertions, explicit
contact-height and scalar inputs, pole obstruction and prepared
realization, and fixed-data simultaneous-prime premises. Its effective
sufficient parameter is $N\ge9682$, with the logarithmic factor in
the contact-height estimate retained. These arithmetic premises are
not independently certified by a finite coefficient verification.

\begin{theorem}\label{thm:main}
Under these premises, the polynomial family specified below consists
of integer quartics $P_N$ in 257 rational variables, including the
input $X$. For every fixed integer $N\ge9682$,
\[
 X\in\Z\quad\Longleftrightarrow\quad
 \exists y_2,\ldots,y_{257}\in\Q:\quad
 P_N(X,y_2,\ldots,y_{257})=0.
\]
For every $N\ge1$ its quadratic specification has 229 rows, its
integer content is one, and its number of nonzero coordinate
monomials is 3679. Its exact coefficient height is
\[
 H(P_N)=\max\{K,149^{2N+2}\},
\]
where
\[
 K=484991541504995159662206970923766531374664106904986910720.
\]
The second term dominates for $N\ge13$. In particular $P_{9682}$
has coefficient height $149^{19366}$.
\end{theorem}

The coefficient specification is
\begin{center}\small
\path{computation/quartic/affine_laurent257_v45/quadratic-family.json},\\
\path{computation/quartic/affine_laurent257_v45/polynomial-family.json}.
\end{center}
The first table lists all quadratic residuals; the second is their
sum of squares, with all coefficients collected. The transformation
certificate in the same directory specifies every coordinate map
and row multiplier used below.

\section{Affine centering of the Lucas residual}

Use the source coordinate numbers, and put
\[
 F=x_{116},\quad F'=x_{117},\quad S=x_{138},\quad S'=x_{139},
 \qquad q=x_{155},\quad w=x_{156}.
\]
The source contains the literal inverse and square equations
\begin{equation}\label{eq:guards}
 qw=1,\qquad x_{157}=w^2.
\end{equation}
Its retained square equations include
\begin{equation}\label{eq:squares}
 F^2+8F'^2=qS,\qquad FF'=qS'.
\end{equation}
Make the triangular affine change
\begin{equation}\label{eq:affine}
 F=\widehat F+q,\qquad
 S=\widehat S+2\widehat F+q,\qquad
 S'=\widehat S'+F',
\end{equation}
fixing every other coordinate. Its matrix has determinant one, and
its inverse has integer coefficients.

Direct substitution turns~\eqref{eq:squares} into
\[
 \widehat F^2+8F'^2=q\widehat S,
 \qquad \widehat F F'=q\widehat S'.
\]
The two expressions controlling the fourth power of the Lucas
residual are
\[
                 S-2F+q,\qquad S'-F'.
\]
Under~\eqref{eq:affine} these become $\widehat S$ and
$\widehat S'$. More explicitly, the source equations for
$G_0=x_{250}$, $G_1=x_{251}$ and $w_0=x_{143}$ become
\[
 w_0G_0=\widehat S^2+32\widehat S'^2,
 \qquad w_0G_1=8\widehat S\widehat S'.
\]
These identities center the residual without introducing a coordinate,
a square witness, or a nonzero condition.

\section{Eliminating an inverse-weight square}

We record the exact equivalence principle used by the construction.

\begin{lemma}\label{lem:laurent}
Suppose a rational polynomial system contains a literal equation
$qw=1$ and a defining equation $z=w^m$. After an invertible affine
change fixing the input, make substitutions
$x_i=w^{a_i}y_i$ with integer exponents, keeping $q,w$ as inverse
coordinates and giving the input exponent zero. Suppose that, after
$z=w^m$ is substituted, each remaining equation becomes a nonzero
Laurent monomial times a polynomial of degree at most two. If the
literal equation $qw=1$ is retained, the new and old rational zero
sets are in input-preserving bijection.
\end{lemma}

\begin{proof}
The retained inverse equation excludes $w=0$ on both zero sets.
Every coordinate substitution and every Laurent row multiplier is
therefore defined and invertible. The deleted coordinate is recovered
uniquely as $w^m$. The affine change is invertible and the input is
fixed. Each direction of the resulting map preserves and reflects
all equations.
\end{proof}

Apply the exponent vector in \path{transformation.json} to the
coordinates centered by~\eqref{eq:affine}. The eliminated coordinate
is $x_{157}=w^2$. The row $qw=1$ is unchanged. The certificate lists,
for each surviving source residual, the integral Laurent exponent and
nonzero rational multiplier relating its substitution to the target
quadratic. Every resulting monomial has coordinate degree at most two.

\begin{proposition}\label{prop:projection}
For every nonzero rational $C$, the certified 257-coordinate system
is in input-preserving bijection with the source rational zero set.
In particular, its sufficient arithmetic parameter remains
$N\ge9682$.
\end{proposition}

\begin{proof}
The exact affine identities are those of~\eqref{eq:affine}. The
Laurent certificate verifies all 229 surviving row identities and
that the deleted square equation becomes an identity. It also
verifies that the literal inverse guard is retained and the input
exponent is zero. Lemma~\ref{lem:laurent} applies. Finally, the
fixed-parameter coordinate scaling described in the next section
is invertible whenever $C\ne0$. Thus the complete zero sets,
and not only a selected collection of witnesses, are equivalent.
The source defining assertion and its effective threshold transfer
without alteration.
\end{proof}

\section{Coefficient scaling and exact height}

The final coordinate change is $y_i=C^{k_i}z_i$, with integer
weights $k_i$ and input weight $k_1=0$. Each quadratic row is divided
by its recorded power of $C$. The resulting rows have integer
coefficients linear in $C$, and the absolute value of every
coefficient multiplying $C$ is at most $149$. This change adds no
coordinate or equation.

Denote these rows by $Q_1,\ldots,Q_{229}$ and define
\[
                    P_N=\sum_{j=1}^{229}Q_j^2.
\]
Every coordinate-monomial coefficient is a polynomial of degree at
most two in $C$. There are 3684 parameter-separated terms, giving
3679 nonzero coordinate-monomial profiles. For each profile $c_m(C)$,
the verifier chooses a sign such that after $C=149+t$ the constant
coefficient is strictly positive and the remaining coefficients are
nonnegative. Therefore none of the profiles vanishes for $C\ge149$.

Both the constant profile $K$ in Theorem~\ref{thm:main} and the
profile $22201C^2$ occur. Exact substitution into every profile gives
height $K$ for $N=1,\ldots,12$. For $C\ge149^{13}$, substituting
$C=149^{13}+t$ into each difference
\[
                   22201C^2\pm c_m(C)
\]
gives only nonnegative coefficients. Hence the height in this range
is exactly $22201C^2=149^{2N+2}$. This proves the stated uniform
height formula.

The input occurs only through its square in the input-square row,
and the coefficient of $X^4$ in the final sum is one. Therefore
$P_N$ has degree exactly four and integer content one for every
allowed parameter. Since a sum of rational squares vanishes precisely
when all summands vanish, Proposition~\ref{prop:projection} proves
Theorem~\ref{thm:main}.

\section{Reproducibility}

The producer uses exact sparse rational arithmetic and a fixed
certificate; it does not invoke a search or optimization routine.
The independent verifier does not import the producer. It checks the
pinned source, all affine and Laurent identities, the deleted
identity, the inverse guard, the fixed input, parameter scaling,
quadratic degrees, and the complete quartic expansion. It verifies
support and height by the exact translated coefficient inequalities
above, uniformly over the parameter range. These are algebraic
certificates and do not replace the retained arithmetic premises.

\begin{thebibliography}{9}
\bibitem{Laurent}
\emph{Quartics in 196 and 258 variables by elimination after monomial
rescaling},
\nolinkurl{papers/quartic/laurent_elimination_20261115_v38.tex}.
\bibitem{Explicit}
\emph{A fully specified defining quartic in 260 variables},
\nolinkurl{papers/quartic/explicit260_20261114_v27.tex}.
\end{thebibliography}
\end{document}
